\begin{alist}
\item
\begin{rlist}
\item Η συνάρτηση $g(x) = f\left(x - \dfrac{\pi}{2}\right)$ προκύπτει από την $f$ αν μετατοπιστεί προς τα δεξιά κατά $\dfrac{\pi}{2}$. Οπότε προκύπτει το παρακάτω σχήμα:
\begin{center}
\begin{tikzpicture}
\begin{axis}[
    axis lines=middle,
    xmin=-1.5*pi/4, xmax=7.5*pi/4,
    ymin=-1.5, ymax=1.5,
    xtick={-pi/4, pi/4, pi/2, 3*pi/4, pi, 5*pi/4, 3*pi/2, 7*pi/4},
    xticklabels={$-\pi/4$, $\pi/4$, $\pi/2$, $3\pi/4$, $\pi$, $5\pi/4$, $3\pi/2$, $7\pi/4$},
    ytick={-1,1},
    grid=both,
    grid style={solid, gray!20},
    thick,
    samples=100
]
\addplot[domain=0:pi, secondarycolor, very thick] {sin(deg(x))};
\node[above] at (axis cs: pi/6, 0.5) {$f$};
\addplot[domain=pi/2:3*pi/2, maincolor, very thick] {sin(deg(x-pi/2))};
\node[above] at (axis cs: pi/2+pi/6+pi/2, 0.5) {$g$};
\end{axis}
\end{tikzpicture}
\end{center}

\item Ο τύπος της $g$ είναι $g(x) = f\left(x - \dfrac{\pi}{2}\right) = \eta\mu\left(x - \dfrac{\pi}{2}\right)$. Η συνάρτηση ορίζεται για $x \in \left[\dfrac{\pi}{2}, \dfrac{3\pi}{2}\right]$, όπως βλέπουμε από το σχήμα.
\end{rlist}

\item Από το α)ii ερώτημα προκύπτει η εξίσωση
\begin{align*}
f(x) = g(x) &\Leftrightarrow \eta\mu x = \eta\mu\left(x - \frac{\pi}{2}\right) \\
&\Leftrightarrow \begin{cases} 
x = 2\kappa\pi + x - \dfrac{\pi}{2} \\ 
x = 2\kappa\pi + \pi - x + \dfrac{\pi}{2} 
\end{cases}, \kappa \in \mathbb{Z} \\
&\Leftrightarrow \begin{cases} 
\text{Αδύνατη} \\ 
2x = 2\kappa\pi + \dfrac{3\pi}{2}, \kappa \in \mathbb{Z} 
\end{cases} \\
&\Leftrightarrow x = \kappa\pi + \dfrac{3\pi}{4}, \kappa \in \mathbb{Z}
\end{align*}
Δεκτή λύση είναι για $\kappa=0$, δηλαδή η γωνία $x = \dfrac{3\pi}{4}$. \\
\textbf{Εναλλακτική λύση:} \\
Η λύση της εξίσωσης $f(x) = g(x)$ είναι η τετμημένη του σημείου τομής των γραφικών παραστάσεων. Οπότε, η μοναδική λύση της εξίσωσης είναι $x = \dfrac{3\pi}{4}$.
\end{alist}
